\documentclass[UTF8, a4paper, 12pt]{article}

\usepackage{fontspec}
\setmainfont{Times New Roman}
\usepackage{xeCJK}
\setCJKmainfont{SimSun}

\usepackage{geometry}
\geometry{margin=1in, top=1.2in, bottom=1.2in}

\usepackage{graphicx}
\usepackage{booktabs}
\usepackage{amsmath, amssymb}

\title{AdvancedRK方法讲解}
\author{讲解人指南}
\date{\today}

\setlength{\parindent}{2em}
\setlength{\parskip}{6pt}

\begin{document}

\maketitle

\section{引言：问题与动机}

\subsection{数据融合的问题}

PM$_{2.5}$监测面临一个基本矛盾：
\begin{itemize}
    \item 地面监测站数据：精度高，但空间覆盖稀疏
    \item CMAQ模型数据：覆盖完整，但存在系统性偏差
\end{itemize}

数据融合的目标：结合两者优点，生成既精确又完整的高分辨率PM$_{2.5}$数据。

\subsection{现有方法的局限}

VNA/eVNA/aVNA方法的核心假设是：CMAQ与真实值之间的偏差是\textbf{线性}关系。

以eVNA为例：假设存在一个常数$r$，使得$O_{true} = O_{CMAQ} \times r$。

这个假设的问题：
\begin{enumerate}
    \item \textbf{非线性偏差}：实际中偏差往往与CMAQ模拟值呈曲线关系（如低浓度高估、高浓度低估）
    \item \textbf{空间异质性}：偏差系数在不同位置可能不同，而不是常数
\end{enumerate}

图1展示了线性假设与实际情况的差异：

\begin{figure}[htbp]
\centering
\fbox{图1：线性假设（虚线）vs 实际非线性偏差（实线）}
\end{figure}

\section{AdvancedRK的核心思想}

AdvancedRK采用\textbf{分步校正}策略：

\textbf{第一步：}用二阶多项式捕捉CMAQ偏差的\textbf{全局非线性趋势}

\textbf{第二步：}用GPR-Matern核捕捉残差的\textbf{局部空间变化}

最终预测：$P = M_{cal} + R^*$，即全局趋势 + 局部调整

\section{第一步详解：二阶多项式趋势校正}

\subsection{为什么需要非线性？}

CMAQ模型的偏差往往不是线性的。原因包括：
\begin{itemize}
    \item 化学反应的非线性：PM$_{2.5}$生成涉及复杂的非线性化学过程
    \item 边界层效应：稳定边界层和混合层对污染物扩散的影响不同
    \item 来源贡献差异：高浓度和低浓度区域的污染来源结构不同
\end{itemize}

用一阶多项式（线性）校正，实质上与eVNA/aVNA没有区别，无法捕捉非线性偏差。

二阶多项式：$M_{cal} = \alpha_0 + \alpha_1 \cdot CMAQ + \alpha_2 \cdot CMAQ^2$

当$\alpha_2 \neq 0$时，表示存在非线性偏差——这是AdvancedRK区别于线性方法的关键。

\subsection{系数的物理含义}

\begin{itemize}
    \item $\alpha_0$：截距项，表示基线偏差（$\mu$g/m$^3$）
    \item $\alpha_1$：一次项系数，若$\alpha_1 > 1$表示CMAQ被低估，$< 1$表示被高估
    \item $\alpha_2$：二次项系数，捕捉非线性程度——当$\alpha_2 \neq 0$时说明偏差呈抛物线模式
\end{itemize}

参数通过普通最小二乘法（OLS）估计，最小化$\sum(O_{monitor} - M_{cal})^2$。

\subsection{为什么选择二阶？}

\begin{itemize}
    \item 一阶（线性）：与VNA/eVNA/aVNA等价，无法捕捉非线性
    \item 二阶：灵活性和复杂度的平衡，能捕捉U形偏差
    \item 三阶及以上：容易过拟合，在新数据上表现差
\end{itemize}

\section{第二步详解：GPR-Matern残差空间建模}

\subsection{为什么要做空间建模？}

多项式校正后，残差$R = O_{monitor} - M_{cal}$应该接近白噪声。但如果残差仍存在空间相关性，说明多项式校正不够——需要用空间插值方法来捕捉这种局部相关性。

图2展示了残差的空间相关性：

\begin{figure}[htbp]
\centering
\fbox{图2：残差的空间分布（相邻点残差相似）}
\end{figure}

\subsection{GPR简介}

高斯过程回归（GPR）是一种非参数的贝叶斯插值方法。

核心思想：假设目标函数（残差）服从高斯过程先验，通过观测数据更新后验分布进行预测。

GPR的优势：
\begin{itemize}
    \item 提供预测值和不确定性估计
    \item 核参数通过最大化边际似然自动确定
    \item 能捕捉复杂的空间相关性结构
\end{itemize}

\subsection{核函数的选择：为什么是Matern？}

核函数决定了"观测点如何影响预测"。

\textbf{RBF核（高斯核）：}
\begin{equation}
k_{RBF}(d) = \sigma^2 \exp\left(-\frac{d^2}{2\ell^2}\right)
\end{equation}

当$d = \ell$时，相关性衰减至$e^{-0.5} \approx 0.37$，之后迅速趋近于0。

\textbf{Matern核（$\nu=1.5$）：}
\begin{equation}
k_{Matern}(d) = \sigma^2 \left(1 + \frac{\sqrt{3}d}{\ell}\right) \exp\left(-\frac{\sqrt{3}d}{\ell}\right)
\end{equation}

图3展示两种核函数的衰减对比：

\begin{figure}[htbp]
\centering
\fbox{图3：Matern核（蓝）vs RBF核（红）衰减对比}
\end{figure}

\textbf{关键区别：}
\begin{itemize}
    \item RBF核：在距离$\ell$处急剧衰减，之后接近0——相当于"突然消失"
    \item Matern核：衰减更平缓，逐渐趋近于0——符合污染物扩散的物理规律
\end{itemize}

\textbf{物理意义：} 想象一个PM$_{2.5}$源区：
\begin{itemize}
    \item 污染物向外扩散时，浓度逐渐降低
    \item 不会在某个距离突然从"有影响"变成"无影响"
    \item 而是随着距离增加，影响逐渐平缓衰减
\end{itemize}

这就是Matern核的物理优势：符合污染物从源区向外扩散的真实过程。

\textbf{参数含义：}
\begin{itemize}
    \item $\ell$（长度尺度）：相关性的特征距离
    \item $\sigma^2$（信号方差）：残差的整体方差规模
    \item $\nu=1.5$（光滑度）：控制函数的光滑程度
\end{itemize}

选择$\nu=1.5$是因为它在"快速衰减"（指数核，$\nu=0.5$）和"极慢衰减"（RBF核，$\nu \to \infty$）之间取得平衡，最符合PM$_{2.5}$扩散规律。

\subsection{克里金预测}

给定新位置$s^*$，残差预测：
\begin{equation}
R^*(s^*) = \mu_R + \mathbf{k}^T \mathbf{K}^{-1} (\mathbf{R} - \mu_R)
\end{equation}

其中：
\begin{itemize}
    \item $\mathbf{k}$：新位置与所有训练点的协方差向量
    \item $\mathbf{K}$：训练点之间的协方差矩阵
    \item $\mathbf{R}$：训练点的残差
\end{itemize}

这本质上是"距离加权平均"的推广——权重由核函数确定。

\section{方法对比与创新点总结}

\begin{table}[htbp]
\centering
\caption{AdvancedRK与基准方法核心差异}
\begin{tabular}{@{}lcc@{}}
\toprule
特征 & 基准方法 & AdvancedRK \\
\midrule
偏差校正 & 线性（常数或简单比例） & \textbf{非线性（二阶多项式）} \\
空间建模 & IDW距离加权 & \textbf{GPR-Matern核} \\
相关性衰减 & 快速（不符合物理） & \textbf{平缓（符合PM$_{2.5}$扩散）} \\
参数估计 & 简单/无训练 & \textbf{OLS+边际似然} \\
不确定性 & 无 & \textbf{有} \\
\end{tabular}
\end{table}

\textbf{核心创新点：}

\begin{enumerate}
    \item \textbf{非线性偏差校正}：二阶多项式能捕捉U形或抛物线形的非线性偏差，线性方法无法做到

    \item \textbf{GPR-Matern核空间建模}：$\nu=1.5$的Matern核衰减更平缓，符合PM$_{2.5}$从源区向外扩散的物理规律

    \item \textbf{两步融合策略}：全局趋势校正 + 局部空间插值，既校正系统性偏差，又保留空间异质性
\end{enumerate}

\section{实验验证}

四个阶段验证结果：

\begin{table}[htbp]
\centering
\caption{AdvancedRK与VNA对比}
\begin{tabular}{@{}lccccccc@{}}
\toprule
 & \multicolumn{3}{c}{AdvancedRK} & \multicolumn{3}{c}{VNA} \\
\cmidrule{2-4} \cmidrule{5-7}
阶段 & $R^2$ & RMSE & |MB| & $R^2$ & RMSE & |MB| & 提升 \\
\midrule
pre\_exp & 0.9047 & 15.58 & 0.02 & 0.8907 & 16.68 & 0.70 & +1.40\% \\
stage1 & 0.9162 & 15.34 & 0.11 & 0.9034 & 16.48 & 0.50 & +1.28\% \\
stage2 & 0.8526 & 4.86 & 0.04 & 0.8408 & 5.05 & 0.05 & +1.18\% \\
stage3 & 0.9129 & 11.57 & 0.02 & 0.9031 & 12.20 & 0.42 & +0.98\% \\
\bottomrule
\end{tabular}
\end{table}

主级创新判定：$R^2 \geq R^2_{VNA} + 0.01$

\textbf{结论：} AdvancedRK在全部四个阶段通过创新判定，R$^2$提升0.98\%到1.40\%，全面优于VNA基准方法。

\section{如何向他人解释}

如果有人问："AdvancedRK是什么？"

简洁回答：
\begin{enumerate}
    \item 用二阶多项式代替线性模型，捕捉CMAQ偏差的非线性关系
    \item 用Matern核代替RBF核，让空间相关性衰减更符合物理规律
    \item 两步融合：全局趋势 + 局部调整
\end{enumerate}

记住核心比喻：
\begin{itemize}
    \item 线性方法 = 用直尺量曲线
    \item AdvancedRK = 用曲线拟合（多项式）+ 高精度空间插值（GPR-Matern）
\end{itemize}

\section{FAQ}

\textbf{Q：为什么要用二阶，三阶不行吗？}

A：二阶是灵活性和复杂度的平衡。三阶以上容易过拟合，在训练数据上表现好但预测新数据时差。就像穿衣服——合身最好，太紧太松都不行。

\textbf{Q：为什么选$\nu=1.5$？}

A：$\nu$控制光滑程度。$\nu=0.5$是指数核（衰减太快），$\nu \to \infty$是RBF核（衰减太慢）。$\nu=1.5$在两者之间，物理上最符合PM$_{2.5}$扩散规律。

\textbf{Q：计算复杂度高吗？}

A：GPR原始复杂度是$O(n^3)$，但我们使用稀疏近似，只保留显著元素（约0.45\%密度），大幅降低计算量同时保持精度。

\end{document}